\chapter{JADWAL DAN RENCANA ANGGARAN BIAYA}
\label{bab:5}

\section{Jadwal Pelaksanaan}
Penelitian tugas akhir ini direncanakan berlangsung secara bertahap selama jangka waktu 4 (empat) bulan aktif. Rincian tahapan kegiatan diuraikan pada Tabel \ref{tab:jadwal} di bawah ini.

\begin{table}[htbp]
\centering
\caption{Jadwal Pelaksanaan Tugas Akhir}
\label{tab:jadwal}
% Mengatur jarak antar kolom menjadi lebih rapat (default biasanya 6pt)
\setlength{\tabcolsep}{3pt} 
\small % Menggunakan font small agar sedikit lebih ringkas namun tetap terbaca jelas
\begin{tabular}{|c|l|c|c|c|c|c|c|c|c|c|c|c|c|c|c|c|c|}
\hline
\multirow{2}{*}{\textbf{No}} & \multicolumn{1}{c|}{\multirow{2}{*}{\textbf{Jenis Kegiatan}}} & \multicolumn{4}{c|}{\textbf{Bulan 1}} & \multicolumn{4}{c|}{\textbf{Bulan 2}} & \multicolumn{4}{c|}{\textbf{Bulan 3}} & \multicolumn{4}{c|}{\textbf{Bulan 4}} \\ \cline{3-18} 
 & \multicolumn{1}{c|}{} & \textbf{1} & \textbf{2} & \textbf{3} & \textbf{4} & \textbf{1} & \textbf{2} & \textbf{3} & \textbf{4} & \textbf{1} & \textbf{2} & \textbf{3} & \textbf{4} & \textbf{1} & \textbf{2} & \textbf{3} & \textbf{4} \\ \hline
1 & Penelusuran pustaka &  $\blacksquare$ & & & & & & & & & & & & & & & \\ \hline
2 & Perancangan rangkaian & & $\blacksquare$ & $\blacksquare$ & & & & & & & & & & & & & \\ \hline
3 & Prototipe & & & & $\blacksquare$ & $\blacksquare$ & & & & & & & & & & & \\ \hline
4 & Instalasi Sistem & & & & & & $\blacksquare$ & $\blacksquare$ & $\blacksquare$ & & & & & & & & \\ \hline
5 & Pengujian & & & & & & & & & $\blacksquare$ & $\blacksquare$ & $\blacksquare$ & $\blacksquare$ & & & & \\ \hline
6 & Laporan Akhir & & & & & & & & & & & & & $\blacksquare$ & $\blacksquare$ & $\blacksquare$ & $\blacksquare$ \\ \hline
\end{tabular}
\end{table}

\section{Rencana Anggaran Biaya}
Rincian anggaran operasional yang dibutuhkan untuk menunjang kebutuhan pembuatan
sistem perangkat keras serta keperluan administrasi dokumen dijabarkan pada
Tabel \ref{tab:rab_alat} dan Tabel \ref{tab:rab_bahan}.

% --- TABEL RAB ALAT ---
\begin{table}[htbp]
\centering
\caption{Rencana Anggaran Biaya - Peralatan Penunjang}
\label{tab:rab_alat}
\begin{tabular}{|c|>{\raggedright\arraybackslash}p{3cm}|>{\raggedright\arraybackslash}p{3.5cm}|c|r|r|} % Kolom 2 dan 3 diatur
\hline
\textbf{No.} & \multicolumn{1}{c|}{\textbf{Material}} & \multicolumn{1}{c|}{\textbf{Justifikasi Pemakaian}} & \textbf{Qty} & \multicolumn{1}{c|}{\textbf{Harga}} & \multicolumn{1}{c|}{\textbf{Total}} \\ \hline
1 & ESP32 Dev Board & Mikrokontroler utama & 1 pcs & 75.000 & 75.000 \\ \hline
2 & IC AM26LS31 & \textit{Differential line driver} & 2 pcs & 12.000 & 24.000 \\ \hline
3 & Optocoupler PC817 & Isolasi galvanik & 4 pcs & 2.500 & 10.000 \\ \hline
4 & Komponen Pasif & Resistor, kapasitor, LED & 1 set & 15.000 & 15.000 \\ \hline
5 & Papan PCB & \textit{Base board} sirkuit & 2 pcs & 20.000 & 40.000 \\ \hline
6 & Kabel \& Jack & Koneksi antarmuka & 1 set & 60.000 & 60.000 \\ \hline
\multicolumn{5}{|r|}{\textbf{Sub Total (Rp.)}} & \textbf{224.000} \\ \hline
\end{tabular}
\end{table}

\newpage
% --- TABEL RAB BAHAN ---
\begin{table}[htbp]
\centering
\caption{Rencana Anggaran Biaya - Bahan Habis Pakai}
\label{tab:rab_bahan}
\begin{tabular}{|c|>{\raggedright\arraybackslash}p{3cm}|>{\raggedright\arraybackslash}p{2.8cm}|c|r|r|} % Kolom 2 dan 3 dipersempit
\hline
\textbf{No.} & \multicolumn{1}{c|}{\textbf{Material}} & \multicolumn{1}{c|}{\textbf{Justifikasi}} & \textbf{Qty} & \multicolumn{1}{c|}{\textbf{Harga}} & \multicolumn{1}{c|}{\textbf{Total}} \\ \hline
1 & Kertas HVS A4 & Laporan & 8 rim & 47.000 & 376.000 \\ \hline
2 & Kertas Art Paper & Poster Laporan & 1 set & 50.000 & 50.000 \\ \hline
3 & Spidol & Tulis menulis & 2 pcs & 7.400 & 14.800 \\ \hline
4 & Pulpen TIZO & Alat tulis & 10 pcs & 5.000 & 50.000 \\ \hline
5 & Kain \& Lakban & Mengemas Alat & 3 pcs & 16.000 & 48.000 \\ \hline
6 & Stabillo Boss & Alat penanda & 1 set & 124.000 & 124.000 \\ \hline
7 & Stiker Transparan & Label Alat & 1 set & 35.000 & 35.000 \\ \hline
8 & Binder Clip & Laporan & 5 set & 11.000 & 55.000 \\ \hline
9 & Pensil 2B & Laporan & 2 pak & 25.000 & 50.000 \\ \hline
\multicolumn{5}{|r|}{\textbf{Sub Total (Rp.)}} & \textbf{802.800} \\ \hline
\multicolumn{5}{|r|}{\textbf{TOTAL KESELURUHAN (Rp.)}} & \textbf{1.026.800} \\ \hline
\end{tabular}
\end{table}